\clearpage\pagestyle{otMain}\setlength{\otBlocEspacement}{2.5mm}\fontsize{11}{15}\selectfont\corrigesChapitre{14}{La droite des milieux}
\begin{blocLivre}{cor-14-1}\begin{corrige}{1}{Une longueur}\demarche{Identifier le troisième côté.}
Dans le triangle ABC, I et J sont les milieux de $[AB]$ et $[AC]$. D’après le théorème du segment des milieux, $IJ=\frac{BC}2=\frac{7{,}4}2=\res{3{,}7\ \mathrm{cm}}$.
\label{fig-cor-14-1}\figureGeo[corrige]{mil-01}
\end{corrige}\end{blocLivre}
\begin{blocLivre}{cor-14-2}\begin{corrige}{2}{Un parallélisme}\demarche{Citer le triangle, les deux milieux et le théorème.}
Dans le triangle EFG, M est le milieu de $[EF]$ et N est le milieu de $[EG]$. D’après le théorème de la droite des milieux, $\res{(MN)\parallel(FG)}$.
\label{fig-cor-14-2}\figureGeo[corrige]{mil-02}
\end{corrige}\end{blocLivre}
\begin{blocLivre}{cor-14-3}\begin{corrige}{3}{Retrouver un côté}\demarche{Doubler la longueur du segment des milieux.}
D’après le théorème du segment des milieux, $IJ=\frac{BC}2$. Donc $BC=2\times5{,}2=\res{10{,}4\ \mathrm{cm}}$.
\label{fig-cor-14-3}\figureGeo[corrige]{mil-03}
\end{corrige}\end{blocLivre}
\begin{blocLivre}{cor-14-4}\begin{corrige}{4}{Un milieu}\demarche{Utiliser un milieu et un parallélisme pour obtenir un milieu.}
Dans le triangle ABC, la droite $(IJ)$ passe par le milieu I de $[AB]$ et elle est parallèle à $(BC)$. D’après le théorème 3, J est le milieu de $[AC]$. Donc $AJ=\frac92=\res{4{,}5\ \mathrm{cm}}$.
\label{fig-cor-14-4}\figureGeo[corrige]{mil-04}
\end{corrige}\end{blocLivre}
\begin{blocLivre}{cor-14-5}\begin{corrige}{5}{Choisir le théorème}\demarche{Repérer ce qui est donné : deux milieux, ou un milieu et un parallélisme.}
a) M et N sont des milieux : théorème 1, $\res{(MN)\parallel(FG)}$. b) M est un milieu et $(MN)\parallel(FG)$ : théorème 3, $\res{N\text{ est le milieu de }[EG]}$. c) M et N sont des milieux : théorème 2, $MN=\frac82=\res{4\ \mathrm{cm}}$.
\label{fig-cor-14-5}\figureGeo[corrige]{mil-05}
\end{corrige}\end{blocLivre}
\begin{blocLivre}{cor-14-6}\begin{corrige}{6}{Le triangle des milieux}\demarche{Associer chaque segment des milieux au côté opposé.}
Chaque côté de IJK joint deux milieux : $IJ=\frac{10}2=5$, $JK=\frac62=3$ et $IK=\frac82=4$ cm. Le périmètre de IJK vaut $\res{12\ \mathrm{cm}}$, soit $\res{\text{la moitié}}$ de celui de ABC, qui vaut 24 cm.
\label{fig-cor-14-6}\figureGeo[corrige]{mil-06}
\end{corrige}\end{blocLivre}
\begin{blocLivre}{cor-14-7}\begin{corrige}{7}{Un milieu, puis une longueur}\demarche{Démontrer d’abord que J est un milieu.}
Étape 1 : $(IJ)$ passe par le milieu I de $[AB]$ et est parallèle à $(BC)$, donc J est le milieu de $[AC]$ (théorème 3). Étape 2 : I et J sont les milieux de $[AB]$ et $[AC]$, donc $IJ=\frac{BC}2=\res{4{,}5\ \mathrm{cm}}$ (théorème 2).
\label{fig-cor-14-7}\figureGeo[corrige]{mil-07}
\end{corrige}\end{blocLivre}
\begin{blocLivre}{cor-14-8}\begin{corrige}{8}{Une rédaction incomplète}\demarche{Vérifier toutes les hypothèses du théorème.}
Le théorème 1 demande deux milieux. Si J n’est pas le milieu de $[AC]$, $(IJ)$ n’est pas parallèle à $(BC)$ : la figure corrigée en donne un exemple. Il faudrait savoir que $\res{J\text{ est le milieu de }[AC]}$.
\label{fig-cor-14-8}\figureGeo[corrige]{mil-08}
\end{corrige}\end{blocLivre}
\begin{blocLivre}{cor-14-9}\begin{corrige}{9}{Dans un rectangle}\demarche{Choisir le triangle qui contient les deux milieux.}
On se place dans le triangle ABC : I et J sont les milieux de $[AB]$ et $[BC]$. D’après les théorèmes 1 et 2, $\res{(IJ)\parallel(AC)}$ et $IJ=\frac{AC}2=\res{5\ \mathrm{cm}}$.
\label{fig-cor-14-9}\figureGeo[corrige]{mil-09}
\end{corrige}\end{blocLivre}
\begin{blocLivre}{cor-14-10}\begin{corrige}{10}{Un parallélogramme de milieux}\demarche{Utiliser deux fois le théorème 1.}
I et K sont les milieux de $[AB]$ et $[BC]$ : $(IK)\parallel(AC)$, donc $(IK)\parallel(AJ)$. J et K sont les milieux de $[AC]$ et $[BC]$ : $(JK)\parallel(AB)$, donc $(JK)\parallel(AI)$. Les côtés opposés de AIKJ sont parallèles deux à deux : c’est un $\res{\text{parallélogramme}}$. Son périmètre vaut $2\times(3{,}5+2{,}5)=\res{12\ \mathrm{cm}}$.
\label{fig-cor-14-10}\figureGeo[corrige]{mil-10}
\end{corrige}\end{blocLivre}
\begin{blocLivre}{cor-14-11}\begin{corrige}{11}{Des milieux de milieux}\demarche{Appliquer le théorème dans deux triangles emboîtés.}
Dans ABC : $IJ=\frac{12}2=6$ cm et $(IJ)\parallel(BC)$. Dans AIJ : M et N sont les milieux de $[AI]$ et $[AJ]$, donc $MN=\frac62=\res{3\ \mathrm{cm}}$ et $(MN)\parallel(IJ)$. Deux droites parallèles à une même troisième sont parallèles : $\res{(MN)\parallel(BC)}$.
\label{fig-cor-14-11}\figureGeo[corrige]{mil-11}
\end{corrige}\end{blocLivre}
\begin{blocLivre}{cor-14-12}\begin{corrige}{12}{Dans un trapèze}\demarche{Utiliser successivement les triangles ABD et BDC.}
Dans ABD, la droite passe par le milieu I de $[AD]$, parallèlement à $(AB)$ : J est le milieu de $[BD]$ et $IJ=\frac62=\res{3\ \mathrm{cm}}$. Dans BDC, elle passe par le milieu J de $[BD]$, parallèlement à $(DC)$ car $(DC)\parallel(AB)$ : K est le milieu de $[BC]$ et $JK=\frac{10}2=\res{5\ \mathrm{cm}}$. Les points sont alignés : $IK=\res{8\ \mathrm{cm}}$.
\label{fig-cor-14-12}\figureGeo[corrige]{mil-12}
\end{corrige}\end{blocLivre}
\begin{blocLivre}{cor-14-13}\begin{corrige}{13}{Le théorème de Varignon}\demarche{Utiliser une diagonale du quadrilatère dans deux triangles.}
Dans ABC : $(IJ)\parallel(AC)$ et $IJ=\frac{AC}2$. Dans ACD : $(LK)\parallel(AC)$ et $LK=\frac{AC}2$. Donc $[IJ]$ et $[LK]$ sont parallèles et de même longueur : IJKL est un $\res{\text{parallélogramme}}$. De même $IL=JK=\frac{BD}2$. Périmètre : $8+6=\res{14\ \mathrm{cm}}$.
\label{fig-cor-14-13}\figureGeo[corrige]{mil-13}
\end{corrige}\end{blocLivre}
\begin{blocLivre}{cor-14-14}\begin{corrige}{14}{Le point de concours des médianes}\demarche{Combiner le théorème des milieux et les diagonales d’un parallélogramme.}
Dans ABG : $(KL)\parallel(AB)$ et $KL=\frac{AB}2$. Dans ABC : $(JI)\parallel(AB)$ et $JI=\frac{AB}2$. Ainsi KLIJ a deux côtés opposés parallèles et de même longueur : c’est un $\res{\text{parallélogramme}}$. Ses diagonales $[KI]$ et $[LJ]$ se coupent en G, leur milieu : $GK=GI$. Comme $AK=KG$, on a $AG=\frac23AI=\res{6\ \mathrm{cm}}$.
\label{fig-cor-14-14}\figureGeo[corrige]{mil-14}
\end{corrige}\end{blocLivre}
\begin{blocLivre}{cor-14-15}\begin{corrige}{15}{Le milieu de l’hypoténuse}\demarche{Transporter l’angle droit par le parallélisme.}
Dans ABC, I et J sont des milieux : $(IJ)\parallel(AC)$. Or $(AC)\perp(AB)$, donc $(IJ)\perp(AB)$. La droite $(IJ)$ est perpendiculaire à $[AB]$ en son milieu J : c’est la $\res{\text{médiatrice de }[AB]}$, donc $IA=IB$. Comme $IB=IC=\frac{BC}2$, on obtient $IA=IB=IC=\res{5\ \mathrm{cm}}$. Ce résultat sera repris au chapitre 16.
\label{fig-cor-14-15}\figureGeo[corrige]{mil-15}
\end{corrige}\end{blocLivre}
